\documentclass[10pt]{article}
\usepackage[utf8]{inputenc}
\usepackage[T1]{fontenc}
\usepackage{lmodern}
\usepackage[margin=1.7cm,top=1.5cm,bottom=1.5cm]{geometry}
\usepackage{amssymb,amsmath}
\usepackage{booktabs}
\usepackage{array}
\usepackage{enumitem}
\usepackage[table]{xcolor}
\usepackage{titlesec}
\usepackage{parskip}

% ---- Paleta de acento (única y consistente) ----
\definecolor{az}{HTML}{1E3A8A}      % azul corporativo (títulos)
\definecolor{azlt}{HTML}{2563EB}    % azul de acento (KPIs clave)
\definecolor{rowbg}{HTML}{EDF1F9}   % sombreado alterno de filas
\definecolor{rulegy}{HTML}{B9C4DA}  % reglas finas / separadores

% ---- Jerarquía tipográfica de secciones ----
\titleformat{\section}
  {\normalfont\sffamily\large\bfseries\color{az}}{}{0em}{}
  [\vspace{1.5pt}{\color{rulegy}\titlerule[0.6pt]}]
\titlespacing{\section}{0pt}{7pt}{3pt}

% ---- Listas: marcadores de color, niveles diferenciados, aire entre ítems ----
\setlist[enumerate,1]{leftmargin=1.7em,itemsep=2.5pt,topsep=2.5pt,parsep=0pt,
  label=\textcolor{az}{\bfseries\arabic*.}}
\setlist[itemize,1]{leftmargin=1.5em,itemsep=2.5pt,topsep=2.5pt,parsep=0pt,
  label=\textcolor{az}{\textbullet}}
\setlist[itemize,2]{leftmargin=1.3em,itemsep=1.5pt,topsep=1.5pt,
  label=\textcolor{azlt}{\small$\triangleright$}}

\setlength{\parskip}{4pt}

% ---- Comandos semánticos de énfasis ----
\newcommand{\stars}{$\star$}
\newcommand{\kpi}[1]{\textbf{\textcolor{azlt}{#1}}}            % número/dato titular
\newcommand{\lead}[1]{\textcolor{az}{\sffamily\bfseries #1}}   % entradilla de párrafo
% ---- Cabecera de anexo consistente con las secciones ----
\newcommand{\anexohead}[1]{%
  \par\vspace{2pt}{\sffamily\normalsize\bfseries\color{az}#1}\par
  \vspace{1.5pt}{\color{rulegy}\hrule height 0.6pt}\vspace{4pt}}
\pagestyle{empty}
\renewcommand{\arraystretch}{1.05}

\begin{document}

\begin{center}
{\LARGE\bfseries\color{az} Informe ejecutivo --- Satisfacción del cliente en Olist}\\[3pt]
{\small\textbf{Foco:} el número de ítems por pedido como palanca de mejora \quad\textbullet\quad
\textbf{Datos:} $\sim$96.000 pedidos entregados (Olist 2016--2018) \quad\textbullet\quad
\textbf{Método:} estadística no paramétrica; decisiones por tamaño de efecto, no por p-valor.}
\end{center}
\vspace{-2pt}\hrule\vspace{4pt}

\section{1. Síntesis}
La satisfacción del cliente (calificación 1--5) se explica, ante todo, por \textbf{cumplir la promesa de entrega}: la demora frente a la fecha prometida es, con diferencia, el factor más importante. Sin embargo, este informe se centra en un segundo hallazgo, menos obvio y accionable de inmediato: \textbf{el número de ítems por pedido}. Es la causa de insatisfacción más fuerte una vez descontado el tiempo, opera por un canal independiente de la logística y hoy pasa desapercibida. Es una palanca barata que conviene controlar antes de que el tamaño promedio de los pedidos crezca.

\section{2. El factor \#1 (contexto): la demora en la entrega}
\begin{itemize}
\item \lead{Es el driver dominante.} Un pedido tardío tiene 5,9$\times$ más riesgo de mala calificación: \kpi{P(calif.\,$\le$2 $\mid$ llega tarde) = 54\%} frente a \textbf{9\%} si llega a tiempo.
\item \lead{El daño es asimétrico:} cada día de retraso resta $\sim$\kpi{0,16 estrellas}, mientras que adelantarse solo suma $\sim$0,01 (asimetría 10,8$\times$). Lo que importa no es la velocidad absoluta, sino \textbf{no cruzar la fecha prometida}.
\item La distancia vendedor--cliente es el motor de esa demora (Spearman \textbf{$\rho=0{,}54$}). \textit{Recomendación permanente: gestionar el ETA y usar ``cumplimiento de la promesa'' como KPI central.}
\end{itemize}

\section{3. Hallazgo central: el número de ítems por pedido}
\lead{Qué encontramos.} A más productos en un mismo pedido, peor tiende a ser la calificación (Spearman \kpi{$\rho=-0{,}107$}). Es la \textbf{segunda causa} de insatisfacción tras el tiempo y la primera de origen no logístico.

\lead{Por qué merece atención --- cuatro razones respaldadas por datos:}
\begin{enumerate}
\item \lead{Es un canal independiente del tiempo de entrega.} El número de ítems no está correlacionado con la demora ($\rho=-0{,}026$) y su efecto sobrevive al controlar por tiempo (correlación parcial $-0{,}117$). No es ``tiempo disfrazado'': es un problema de \textbf{cumplimiento del pedido}, no de velocidad, atacable por separado.
\item \lead{El efecto condicional es grande, aunque el promedio lo diluya.} En pedidos entregados a tiempo (para aislar el efecto), la tasa de clientes muy molestos (1--2\stars) pasa de \kpi{7,5\% con 1 ítem} a \kpi{23\% con 2 ítems} (¡el triple!) y \kpi{32\% con 5+}. El $\rho$ global parece pequeño ($\sim$1\% de la varianza) solo porque hoy $\sim$90\% de los pedidos son de un único ítem.
\item \lead{Mecanismo: ``eslabón más débil''.} Cada ítem adicional es otra oportunidad de que algo falle (dañado, incompleto, empaque distinto, entrega parcial); basta que un producto falle para hundir la nota de todo el pedido. La evidencia lo confirma: cuando el pedido llega de varios vendedores (paquetes separados), el descontento salta de \kpi{21\% a 47\%}.
\item \lead{Es real, no azar, y su magnitud es sólida (Monte Carlo).} Una prueba de permutación (5.000 simulaciones que rompen el vínculo ítems--insatisfacción) produjo una brecha máxima por azar de \textbf{1,05 pp}; la brecha real es de \kpi{17,1 pp} (\textbf{p\,$<$\,0,0002}). Un bootstrap (4.000 remuestreos) estima $\sim$\textbf{173 clientes insatisfechos evitables por cada 10.000 pedidos} hoy, con IC 90\% [165--181] (y [156--190] con \emph{bootstrap por vendedor}, más conservador).
\end{enumerate}

\lead{Por qué actuar ahora (proyección).} Hoy el número de ítems explica solo $\sim$\kpi{1,1\% de la varianza} de la satisfacción ($\rho^2$ con $\rho=-0{,}107$), un efecto pequeño \emph{únicamente} porque $\sim$90\% de los pedidos son de un ítem; \textbf{no es una cifra fija}. Según las simulaciones, el impacto escala linealmente con la proporción de pedidos multi-ítem: si esa proporción crece (más venta cruzada, canastas más grandes), los detractores evitables suben de $\sim$173 hoy a \kpi{$\sim$513 (si el multi-ítem llega al 30\%)} y \kpi{$\sim$855 (al 50\%)} por cada 10.000 pedidos. Es barato prevenirlo ahora; se encarece si se deja crecer.

\section{4. Comportamiento del vendedor (mención breve)}
El desempeño operativo del vendedor lo explica sobre todo su \textbf{volumen} ($\rho=0{,}20$ con el \% de despachos tardíos; distancia y región son irrelevantes). Clave para la interpretación: \kpi{el vendedor solo controla $\sim$11--14\% de su propia calificación}; el resto depende del transportador y del producto. No debe penalizarse al vendedor por demoras del transportista.

\section{5. Recomendaciones accionables}
\begin{enumerate}
\item \lead{Programa de cumplimiento para pedidos multi-ítem} (prioridad de este informe): verificación de completitud antes del despacho, empaque reforzado y \textbf{consolidación en un solo envío} cuando sea posible. Independiente de la logística de entrega y de bajo costo.
\item \lead{Marcar como ``alto riesgo'' los pedidos multi-ítem y multi-vendedor} para control de calidad proactivo; son los del efecto ``eslabón más débil''.
\item \lead{Mantener la entrega como prioridad \#1:} gestión de ETA, KPI de cumplimiento de promesa y auditoría logística automática ante cada 1\stars\ (una 1\stars\ es 4,7$\times$ más probable si hubo retraso).
\item \lead{Vendedor:} tutelar a los sellers según volumen y usar el \textbf{\% de despacho tardío} como KPI atribuible; no castigarlos por el transportador.
\item \lead{No invertir} en método de pago, flete ni categoría como palancas de satisfacción: su efecto práctico es nulo (ver Anexo).
\end{enumerate}

\vspace{2pt}\hrule\vspace{3pt}
\anexohead{Anexo A --- Evidencia cuantitativa de los 11 conceptos}
{\footnotesize
\renewcommand{\arraystretch}{1.25}
\rowcolors{2}{rowbg}{white}
\begin{tabular}{@{}>{\centering\arraybackslash}m{1.1em} l l@{}}
\toprule
\rowcolor{white}\textbf{\#} & \textbf{Concepto} & \textbf{Resultado} \\
\midrule
1 & Probabilidad condicional & P($\le$2\stars$\mid$tarde)=0,54 vs 0,09 a tiempo (RR 5,9$\times$) \\
2 & Teorema de Bayes & P(tarde$\mid$1\stars)=0,38 (prior 0,08); lift 4,7$\times$ \\
3 & Verosimilitud / MLE & Flete: log-normal ($\sigma$=0,52; escala 17,33); mejor AIC \\
4 & Distribuciones paramétricas & Tiempo de entrega y precio: log-normales (KS 0,032 y 0,029) \\
5 & Esperanza y varianza & E=4,09; Var=1,82; peores estados AL/MA $\approx$3,76 \\
6 & Independencia y correlación & tiempo $-0{,}235$; retraso $-0{,}176$; \textbf{núm.\ ítems $-0{,}107$}; flete $-0{,}088$; distancia $-0{,}065$ \\
7 & Prior y posterior & Shrinkage: $\kappa$=1,92; seller 1 pedido/5,0\stars\ $\to$ posterior 4,33 \\
8 & Entropía (Shannon) & pago 1,02 bits (efic.\ 0,44); categoría 4,73 bits (0,77) \\
9 & Entropía cruzada & Clasificador reseña: log-loss 0,471 vs base 0,515 ($-8{,}5\%$); AUC 0,66 \\
10 & Divergencia KL & a-tiempo$\|$tarde 0,81 bits; SP$\|$Nordeste 0,03 bits \\
11 & $\chi^2$ (independencia) & Pago $\times$ reseña: $\chi^2$=32,4; \textbf{V de Cramér=0,018} (efecto nulo) \\
\bottomrule
\end{tabular}}

\vspace{6pt}
\anexohead{Anexo B --- Cifras del hallazgo ``número de ítems'' (pedidos entregados a tiempo)}
{\footnotesize
\renewcommand{\arraystretch}{1.25}
\rowcolors{2}{rowbg}{white}
\begin{tabular}{@{}l l@{}}
\toprule
\rowcolor{white}\textbf{Métrica} & \textbf{Valor} \\
\midrule
\% detractores (1--2\stars) por n.\ de ítems & 1: 7,5\% \ $\cdot$\ 2: 23\% \ $\cdot$\ 3: 27\% \ $\cdot$\ 4: 32\% \ $\cdot$\ 5+: 32\% \\
\% detractores en multi-ítem según vendedores & 1 vendedor 21\% \ $\cdot$\ varios 47\% \\
Correlación n.\ ítems $\leftrightarrow$ tiempo de entrega & $\rho=-0{,}026$ (independiente del tiempo) \\
Correlación parcial (controlando tiempo) & $-0{,}117$ (el efecto persiste) \\
Monte Carlo --- permutación & brecha real 17,1 pp vs máx.\ azar 1,05 pp; p\,$<$\,0,0002 \\
Bootstrap --- evitables/10.000 pedidos (hoy) & 173 (IC 90\%: 165--181; por vendedor 156--190) \\
Proyección si multi-ítem sube a 30\% / 50\% & $\sim$513 / $\sim$855 evitables por 10.000 pedidos \\
\bottomrule
\end{tabular}}

\vspace{4pt}
{\scriptsize\textit{Nota metodológica: asociación observacional (no causal); datos 2016--2018. Detalles reproducibles en \texttt{outputs/reports/}; dashboard interactivo en \texttt{dashboard/index.html}.}}

\end{document}
